% SPDX-License-Identifier: PMPL-1.0-or-later
% Copyright (c) 2026 Jonathan D.A. Jewell
%
% Dual-Form Language Design: Turing-Incomplete Constrained Deployment
% with Turing-Complete Factory Compilation
%
% arXiv-style preprint

\documentclass[11pt,a4paper]{article}

% --- Packages ---
\usepackage[utf8]{inputenc}
\usepackage[T1]{fontenc}
\usepackage{lmodern}
\usepackage{amsmath,amssymb,amsthm}
\usepackage{stmaryrd}
\usepackage{mathtools}
\usepackage{microtype}
\usepackage{hyperref}
\usepackage{cleveref}
\usepackage{listings}
\usepackage{xcolor}
\usepackage{booktabs}
\usepackage{enumitem}
\usepackage{tikz}
\usetikzlibrary{arrows.meta,positioning,shapes.geometric,calc,fit,backgrounds}
\usepackage{algorithm}
\usepackage{algpseudocode}
\usepackage[margin=1in]{geometry}
\usepackage{graphicx}
\usepackage{natbib}
\bibliographystyle{plainnat}
\usepackage{doi}

% --- Theorem environments ---
\newtheorem{theorem}{Theorem}[section]
\newtheorem{lemma}[theorem]{Lemma}
\newtheorem{proposition}[theorem]{Proposition}
\newtheorem{corollary}[theorem]{Corollary}
\newtheorem{definition}[theorem]{Definition}
\newtheorem{remark}[theorem]{Remark}

% --- Listings configuration ---
\definecolor{oblkw}{RGB}{0,0,160}
\definecolor{oblcmt}{RGB}{80,120,80}
\definecolor{oblstr}{RGB}{160,32,32}
\definecolor{oblrev}{RGB}{140,60,180}
\definecolor{obltrace}{RGB}{0,120,120}
\definecolor{oblfact}{RGB}{180,100,0}

\lstdefinelanguage{oblibeny}{
  morekeywords={fn,let,mut,for,in,if,else,match,return,true,false,const,struct,mod},
  morekeywords=[2]{swap,incr,decr,push,pop,xor_assign},
  morekeywords=[3]{trace,checkpoint,assert_invariant},
  morekeywords=[4]{define,emit,quote_cf,gensym,splice,when_emit,lambda,letrec},
  sensitive=true,
  morecomment=[l]{//},
  morecomment=[s]{(*}{*)},
  morestring=[b]",
  keywordstyle=\color{oblkw}\bfseries,
  keywordstyle=[2]\color{oblrev}\bfseries,
  keywordstyle=[3]\color{obltrace}\bfseries,
  keywordstyle=[4]\color{oblfact}\bfseries,
  commentstyle=\color{oblcmt}\itshape,
  stringstyle=\color{oblstr},
  basicstyle=\ttfamily\small,
  breaklines=true,
  showstringspaces=false,
  tabsize=2,
  frame=single,
  framerule=0.4pt,
  rulecolor=\color{gray!40},
  backgroundcolor=\color{gray!3},
}

\lstdefinelanguage{ocaml-ast}{
  morekeywords={type,and,of,match,with,let,rec,fun,function,if,then,else,
                module,struct,end,sig,val,open,Some,None,true,false,int,string,
                bool,list,option,ref,mutable},
  sensitive=true,
  morecomment=[s]{(*}{*)},
  morestring=[b]",
  keywordstyle=\color{oblkw}\bfseries,
  commentstyle=\color{oblcmt}\itshape,
  stringstyle=\color{oblstr},
  basicstyle=\ttfamily\small,
  breaklines=true,
  showstringspaces=false,
  tabsize=2,
  frame=single,
  framerule=0.4pt,
  rulecolor=\color{gray!40},
  backgroundcolor=\color{gray!3},
}

\lstset{
  defaultdialect=[]{oblibeny},
  basicstyle=\ttfamily\small,
  breaklines=true,
}

% --- Custom commands ---
\newcommand{\obl}{\textsc{Obl\'{i}ben\'{y}}}
\newcommand{\cfform}{\textsf{CF}}
\newcommand{\ffform}{\textsf{FF}}
\newcommand{\bigO}{\mathcal{O}}
\newcommand{\sem}[1]{\llbracket #1 \rrbracket}
\newcommand{\inv}[1]{#1^{-1}}

% --- Title ---
\title{%
  Dual-Form Language Design:\\
  Turing-Incomplete Constrained Deployment\\
  with Turing-Complete Factory Compilation%
}

\author{%
  Jonathan D.A.\ Jewell\\
  \texttt{j.d.a.jewell@open.ac.uk}\\
  \texttt{https://github.com/hyperpolymath}
}

\date{March 2026}

\begin{document}

\maketitle

% ============================================================================
% ABSTRACT
% ============================================================================
\begin{abstract}
We present \obl{}, a dual-form programming language in which a single source
language separates into two distinct computational regimes: a
\emph{Turing-complete factory form} used at compile time for metaprogramming,
code generation, and optimisation, and a \emph{Turing-incomplete constrained
form} that constitutes the deployed runtime artifact.  The constrained form
enforces termination, bounded resource consumption, full reversibility of all
state transitions, and cryptographically-chained accountability traces---by
\emph{syntactic and structural restriction}, not by optional annotation or
runtime monitoring.  We formalise the dual-form compilation model, prove that
constrained-form programs terminate within statically-computable resource
bounds (the \emph{Maximum Principle Reduction}), and demonstrate that all
constrained-form operations are invertible with respect to a well-defined
inverse semantics.  The compiler is implemented in OCaml with formal
interface proofs in Idris2 and a post-quantum cryptographic FFI via Zig.  We
evaluate the design against deployment requirements for Hardware Security
Modules, secure enclaves, and critical embedded systems, where
non-termination, unauditable state, and irreversible failure modes represent
not merely inconveniences but regulatory and safety violations.
\end{abstract}

\medskip
\noindent\textbf{Keywords:} Turing-incompleteness, total functional programming,
reversible computing, dual-form compilation, hardware security modules,
post-quantum cryptography, formal verification, accountability

% ============================================================================
% 1. INTRODUCTION
% ============================================================================
\section{Introduction}
\label{sec:introduction}

The overwhelming majority of deployed software is written in Turing-complete
languages.  This is simultaneously the discipline's greatest achievement and
its most persistent liability.  Turing-completeness grants unbounded
expressiveness: any computable function can be expressed.  But this
expressiveness comes at a price that is paid in every safety-critical,
security-critical, and compliance-critical deployment: \emph{the halting
problem is undecidable}, and therefore no general algorithm can determine
whether an arbitrary deployed program will terminate, consume bounded
resources, or avoid catastrophic state corruption.

For the majority of application software---web services, desktop applications,
batch processing---the consequences of non-termination are manageable.
Watchdog timers, process supervisors, and human operators provide external
termination guarantees.  But there exists a class of deployment targets where
these external mechanisms are either unavailable, unacceptable, or
insufficient:

\begin{itemize}[nosep]
  \item \textbf{Hardware Security Modules (HSMs):}  Firmware executing inside
    tamper-resistant cryptographic processors, where an infinite loop can brick
    a device that costs tens of thousands of dollars and may be physically
    inaccessible.
  \item \textbf{Secure enclaves} (Intel SGX, ARM TrustZone):  Code executing
    in isolated memory regions where external monitoring is architecturally
    prohibited---the enclave \emph{cannot} be interrupted by the host OS
    without destroying its security guarantees.
  \item \textbf{Smart cards and embedded tokens:}  Devices with no watchdog
    capability, no operating system, and no mechanism for external recovery
    from hung execution.
  \item \textbf{Critical embedded systems:}  Automotive, aerospace, and
    medical device firmware where non-termination of a control loop can cause
    physical harm.
  \item \textbf{Regulatory environments:}  Financial, healthcare, and
    government systems where every state transition must be auditable, every
    operation reversible for dispute resolution, and every computation
    provably terminating.
\end{itemize}

The standard response to these concerns has been to use Turing-complete
languages and then \emph{verify} that specific programs terminate---via model
checking, abstract interpretation, or theorem proving applied post-hoc.  This
approach works, but it is expensive, fragile, and places the verification
burden on the \emph{user} of the language rather than on its \emph{designer}.
Every new program requires new verification effort.  Every modification
requires re-verification.

We propose a different approach: \textbf{make non-termination
\emph{inexpressible}}.

\obl{} (Czech: ``beloved'') is a language in which the deployed runtime
artifact is \emph{syntactically incapable} of non-termination.  Not because a
clever verifier has proved termination of a particular program, but because
the language's grammar and type system physically lack the constructs needed
to express unbounded computation.  There is no \texttt{while}.  There is no
\texttt{loop}.  There is no recursion---direct or indirect.  The only
iteration construct is \texttt{for}~$i$~\texttt{in}~$0$\texttt{..}$N$, where
$N$ must be a compile-time constant, and the call graph is statically verified
to be acyclic.

But we do not stop at termination.  The constrained form also enforces:

\begin{enumerate}[nosep]
  \item \textbf{Bounded resource consumption:}  Every program's maximum
    execution time and memory usage are statically computable before execution
    begins.
  \item \textbf{Full reversibility:}  Every state transition has a
    well-defined inverse.  Any execution can be run backward to reconstruct
    prior states without explicit undo logic.
  \item \textbf{Complete accountability:}  Every operation appends to an
    immutable, cryptographically-chained trace sufficient to replay, audit, or
    reverse the entire computation.
\end{enumerate}

The obvious objection is that Turing-incompleteness is too restrictive.  How
can a useful program be written without general recursion or unbounded loops?
Our answer is the \emph{dual-form architecture}: developers write in a
Turing-complete \emph{factory form} that has full metaprogramming
capabilities.  The factory form executes at compile time and \emph{emits}
constrained-form programs as its output.  The developer retains full
expressive power; the deployed artifact retains full safety guarantees.

The key insight is that \textbf{you do not need Turing-completeness at
runtime} for the vast majority of deployed systems.  The factory produces the
constrained artifact, and the constrained artifact is what faces the hostile
environment.

This paper makes the following contributions:

\begin{enumerate}[nosep]
  \item We formalise the dual-form compilation model with a precise definition
    of the factory-to-constrained form boundary (\cref{sec:dual-form}).
  \item We prove the \emph{Maximum Principle Reduction}: every constrained-form
    program has a statically-computable upper bound on both execution time and
    memory consumption (\cref{sec:maximum-principle}).
  \item We define a reversibility semantics for the constrained form and prove
    that all operations are invertible with respect to this semantics
    (\cref{sec:reversibility}).
  \item We describe the OCaml-based compiler implementation and its integration
    with Idris2 formal proofs and Zig-based post-quantum FFI
    (\cref{sec:implementation}).
  \item We evaluate the design against real deployment scenarios in HSM
    firmware and secure enclave applications (\cref{sec:case-studies}).
\end{enumerate}


% ============================================================================
% 2. BACKGROUND
% ============================================================================
\section{Background and Related Work}
\label{sec:background}

\subsection{Total Functional Programming and Turing-Incompleteness}

The idea of restricting computation to guarantee termination has a long
lineage.  Turner's \emph{Total Functional Programming}
\citep{turner2004total} proposed that practical programming could be done in a
language where every function is total---that is, defined on all inputs and
guaranteed to produce a result.  Turner identified the key technical
mechanisms: primitive recursion over algebraic data types (where every
recursive call operates on a structurally-smaller argument) and codata for
potentially-infinite structures.

The \emph{Charity} language \citep{cockett1992charity} demonstrated that
categorical semantics could enforce totality: Charity programs correspond to
morphisms in distributive categories, and the categorical structure prevents
partiality.  The \emph{Walrus} language explored similar ideas in a more
imperative setting.

Modern dependently-typed languages have made totality checking practical:
\begin{itemize}[nosep]
  \item \textbf{Agda} \citep{norell2007towards} includes a termination
    checker based on structural recursion with size types, rejecting programs
    that cannot be shown to terminate.
  \item \textbf{Idris2} \citep{brady2021idris2} uses quantitative type
    theory to track linearity and totality, with an optional totality
    annotation system.
  \item \textbf{Coq}'s \citep{coq2024} guard condition ensures that all
    fixpoints are structurally decreasing.
\end{itemize}

\obl{} differs from these approaches in a fundamental way: totality is not
enforced by a \emph{termination checker} that analyses recursion patterns.  It
is enforced by the \emph{absence} of recursion from the language entirely.
This is a stronger guarantee---there is no termination checker to be confused
by complex recursion schemes, because there are no recursion schemes.  The
trade-off is expressiveness, which we recover through the factory form.

\subsection{Reversible Computing}

Reversible computing---computation in which every step can be undone---has
been studied from both thermodynamic and computational perspectives.
Landauer's principle \citep{landauer1961irreversibility} established that
logically irreversible operations dissipate a minimum amount of energy,
motivating the study of logically reversible computation as a path toward
thermodynamically efficient computing.

The \emph{Janus} language \citep{yokoyama2007reversible} is the canonical
reversible imperative language: every Janus statement has a unique inverse,
and programs can be run forward or backward.  Janus achieves this through
careful restriction of assignment (to reversible operations like XOR-swap and
increment/decrement) and paired control structures (if-then-else with
assertions that determine the entry branch during backward execution).

\emph{R-WHILE} \citep{gluck2012rwhile} and \emph{R-CORE} extend reversible
programming to while-loops with reversible step-operators.

\obl{}'s reversibility differs from Janus in two respects.  First, \obl{}'s
reversibility operates within a Turing-incomplete fragment, meaning that
backward execution is also guaranteed to terminate---a property that Janus,
being Turing-complete, cannot guarantee in general.  Second, \obl{} uses an
\emph{accountability trace} to record information that would otherwise be
destroyed, enabling operations like overwrite (which is irreversible in
isolation) to become reversible in the presence of the trace.  This is closer
to Bennett's \emph{reversible pebble game} \citep{bennett1989timespace} than
to Janus's strict logical reversibility.

\subsection{Hardware Security Modules and Secure Enclaves}

HSMs are dedicated cryptographic processors that manage digital keys and
perform encryption/decryption operations within a physically-protected
boundary.  The FIPS~140-3 standard \citep{fips1403} and Common Criteria
\citep{commoncriteria} impose strict requirements on HSM firmware:
deterministic execution, bounded response times, complete audit trails, and
the ability to roll back to known-good states.

Intel SGX \citep{costan2016intel} and ARM TrustZone \citep{pinto2019demystifying}
provide software-based secure enclaves with similar
properties: isolated execution, attestation, and sealed storage.  Code running
inside an enclave cannot be debugged or monitored by the host operating
system, which means that traditional watchdog-based termination guarantees
are unavailable.

Current practice for HSM and enclave firmware development relies on
Turing-complete languages (typically C or C++) with extensive manual review,
static analysis, and sometimes formal verification of specific components.
The cost of this approach is measured in person-years per firmware release.
\obl{} aims to shift the guarantee from verification of specific programs to
a \emph{language-level property} that holds for all programs.

\subsection{Post-Quantum Cryptography}

The NIST Post-Quantum Cryptography standardisation process
\citep{nist2024pqc} has selected ML-KEM (Kyber) for key encapsulation and
ML-DSA (Dilithium) and SLH-DSA (SPHINCS+) for digital signatures.  HSM
firmware must incorporate these algorithms as classical algorithms face
deprecation.  \obl{}'s foreign function interface provides access to these
primitives through a Zig-based binding to liboqs \citep{stebila2024liboqs}
and libsodium, with Idris2-verified ABI definitions ensuring type-level
correctness of the binding.


% ============================================================================
% 3. THE DUAL-FORM MODEL
% ============================================================================
\section{The Dual-Form Model}
\label{sec:dual-form}

\subsection{Overview}

\obl{} defines two syntactic forms that share a common type system but differ
in computational power:

\begin{definition}[Factory Form \ffform{}]
The \emph{factory form} is a Turing-complete expression language with
S-expression syntax, supporting lambda abstraction, recursive binding
(\texttt{letrec}), pattern matching, and code generation primitives
(\texttt{emit}, \texttt{quote-cf}, \texttt{gensym}, \texttt{splice}).
Factory-form programs execute at compile time and produce constrained-form
programs as output.
\end{definition}

\begin{definition}[Constrained Form \cfform{}]
The \emph{constrained form} is a Turing-incomplete imperative language with
C-like syntax, supporting fixed-width integer types, immutable and mutable
bindings, bounded iteration, acyclic function calls, reversible primitives,
and accountability trace operations.  The constrained form is the only form
that can execute at runtime.
\end{definition}

The compilation pipeline is:

\begin{center}
\begin{tikzpicture}[
  box/.style={draw, rounded corners, minimum width=2.8cm, minimum height=1cm,
              align=center, font=\small},
  arr/.style={-{Stealth[length=5pt]}, thick},
  lbl/.style={font=\footnotesize\itshape, midway, above}
]
  \node[box, fill=orange!15] (src) {\ffform{} Source\\(Turing-complete)};
  \node[box, fill=blue!10, right=2.5cm of src] (cf) {\cfform{} Program\\(Turing-incomplete)};
  \node[box, fill=green!10, right=2.5cm of cf] (bin) {Deployed\\Binary};

  \draw[arr] (src) -- node[lbl] {factory eval} (cf);
  \draw[arr] (cf) -- node[lbl] {constrained compile} (bin);

  \node[below=0.3cm of src, font=\footnotesize\color{gray}] {Development time};
  \node[below=0.3cm of cf, font=\footnotesize\color{gray}] {Verified artifact};
  \node[below=0.3cm of bin, font=\footnotesize\color{gray}] {Runtime};
\end{tikzpicture}
\end{center}

\subsection{Factory Form Syntax and Semantics}

The factory form uses S-expression syntax inspired by Scheme, with extensions
for constrained-form generation:

\begin{lstlisting}[language=oblibeny,caption={Factory form generating a bounded counter}]
(define (make-bounded-counter max-val)
  (emit
    `(fn counter() -> i32 {
       let mut c: i32 = 0;
       for i in 0..,(quote_cf max-val) {
         incr(c, 1);
         trace("increment", c);
       }
       c
     })))
\end{lstlisting}

Invoking \texttt{(make-bounded-counter 100)} at factory time produces a
constrained-form function with a statically-bounded loop of exactly 100
iterations.  The factory form can compute the bound dynamically---it is
Turing-complete---but the emitted artifact has a fixed bound.

\begin{definition}[Factory-Form Well-Formedness]
A factory-form program $P_F$ is well-formed if:
\begin{enumerate}[nosep,label=(F\arabic*)]
  \item\label{wf:f1} Every expression passed to \texttt{emit} is a syntactically valid
    constrained-form expression.
  \item\label{wf:f2} Factory-form execution is bounded by a configurable resource limit
    (time and memory).  If the limit is exceeded, compilation fails.
  \item\label{wf:f3} Factory-form expressions cannot reference runtime values; the
    factory and constrained forms share no mutable state.
\end{enumerate}
\end{definition}

Rule~\ref{wf:f2} deserves comment: we do not solve the halting problem for
factory-form programs.  Instead, we impose a \emph{practical} resource bound
on factory-time execution, analogous to template instantiation depth limits in
C++ or macro expansion limits in Rust.  If a factory program does not
terminate within the bound, compilation is aborted.  This is acceptable
because the factory executes in a development environment where retry and
debugging are available.

\subsection{Constrained Form Syntax}

The constrained form uses imperative syntax with explicit type annotations:

\begin{lstlisting}[language=oblibeny,caption={Constrained form---complete valid program}]
fn main() -> () {
  let mut a: i64 = 0;
  let mut b: i64 = 1;
  checkpoint("start_fibonacci");
  for i in 0..10 {
    trace("fib_value", a);
    let tmp: i64 = a;
    swap(a, b);
    incr(b, tmp);
  }
  checkpoint("end_fibonacci");
  assert_invariant(a == 55, "10th Fibonacci is 55");
}
\end{lstlisting}

The constrained form's syntax is defined by the following grammar (simplified):

\begin{align*}
  \textit{prog}  &::= \textit{decl}^* \\
  \textit{decl}  &::= \texttt{fn}\; \textit{name}\texttt{(}\textit{params}\texttt{)} \to \textit{type}\; \texttt{\{}\; \textit{stmt}^*\; \texttt{\}} \\
                 &\;|\; \texttt{struct}\; \textit{name}\; \texttt{\{}\; (\textit{name}\texttt{:}\;\textit{type})^*\; \texttt{\}} \\
                 &\;|\; \texttt{const}\; \textit{name}\texttt{:}\;\textit{type}\; \texttt{=}\; \textit{expr} \\[6pt]
  \textit{stmt}  &::= \texttt{let}\; \textit{name}\texttt{:}\;\textit{type}\; \texttt{=}\; \textit{expr}
                 \;|\; \texttt{let\;mut}\; \textit{name}\texttt{:}\;\textit{type}\; \texttt{=}\; \textit{expr} \\
                 &\;|\; \textit{name}\; \texttt{=}\; \textit{expr}
                 \;|\; \texttt{if}\; \textit{expr}\; \texttt{\{}\;\textit{stmt}^*\;\texttt{\}}\; \texttt{else}\; \texttt{\{}\;\textit{stmt}^*\;\texttt{\}} \\
                 &\;|\; \texttt{for}\; \textit{name}\; \texttt{in}\; n\texttt{..}N\; \texttt{\{}\;\textit{stmt}^*\;\texttt{\}}
                       \quad \text{where } n, N \in \mathbb{Z}_{\text{const}} \\
                 &\;|\; \texttt{swap}\texttt{(}\textit{a}\texttt{,}\;\textit{b}\texttt{)}
                 \;|\; \texttt{incr}\texttt{(}\textit{x}\texttt{,}\;\textit{e}\texttt{)}
                 \;|\; \texttt{decr}\texttt{(}\textit{x}\texttt{,}\;\textit{e}\texttt{)} \\
                 &\;|\; \textit{x}\; \hat{}\texttt{=}\; \textit{e}
                 \;|\; \texttt{trace}\texttt{(}\textit{tag}\texttt{,}\;\textit{args}\texttt{)}
                 \;|\; \texttt{checkpoint}\texttt{(}\textit{label}\texttt{)} \\
                 &\;|\; \texttt{assert\_invariant}\texttt{(}\textit{expr}\texttt{,}\;\textit{msg}\texttt{)}
\end{align*}

Crucially, the grammar does not contain:
\begin{itemize}[nosep]
  \item \texttt{while} or \texttt{loop} --- there is no production rule for
    unbounded iteration.
  \item Self-referential function definitions --- no production permits a
    function body to reference its own name in a call position.
  \item Indirect recursion is prevented not by the grammar but by a
    \emph{call-graph acyclicity check} that runs as part of compilation.
\end{itemize}

\subsection{Formal Relationship Between Forms}

We can characterise the relationship between the two forms precisely.

\begin{definition}[Form Embedding]
Let $\mathcal{L}_F$ be the set of all syntactically-valid factory-form
programs and $\mathcal{L}_C$ be the set of all syntactically-valid
constrained-form programs.  The factory evaluation function
\[
  \mathsf{eval}_F : \mathcal{L}_F \rightharpoonup \mathcal{L}_C
\]
is a partial function (partial because factory evaluation may diverge or
exceed resource bounds) that maps factory programs to constrained programs.
\end{definition}

\begin{proposition}[Constrained Closure]
For any factory program $P_F \in \mathcal{L}_F$, if
$\mathsf{eval}_F(P_F) = P_C$, then $P_C \in \mathcal{L}_C$.  That is, the
factory can only produce valid constrained-form programs.
\end{proposition}

\begin{proof}
By rule~\ref{wf:f1}: the \texttt{emit} construct accepts only
syntactically-valid constrained-form expressions, enforced by the factory
evaluator's type system.  All emitted fragments are parsed and validated
against the constrained-form grammar before assembly.
\end{proof}

\begin{remark}[The Factory as a Staging System]
The dual-form model is a form of \emph{multi-stage programming}
\citep{taha2004gentle}, where the factory form is stage~0 (compile time) and
the constrained form is stage~1 (runtime).  The key difference from
MetaOCaml-style staging is that the stages have \emph{different computational
power}: stage~0 is Turing-complete, stage~1 is Turing-incomplete.  This
asymmetric staging is, to our knowledge, novel.
\end{remark}


% ============================================================================
% 4. MAXIMUM PRINCIPLE REDUCTION
% ============================================================================
\section{Maximum Principle Reduction}
\label{sec:maximum-principle}

The central guarantee of the constrained form is not merely termination, but
\emph{bounded termination}: every program's maximum execution time and memory
consumption can be computed statically, before execution begins.  We call this
the \emph{Maximum Principle Reduction} (MPR), by analogy with the maximum
principle in PDE theory---just as the maximum of a harmonic function occurs on
the boundary of its domain, the maximum resource consumption of a
constrained-form program is determined at its ``boundary'' (the compile-time
constants that parameterise its loops).

\subsection{Resource Model}

\begin{definition}[Resource Certificate]
A \emph{resource certificate} for a constrained-form program $P_C$ is a tuple
$\mathcal{R}(P_C) = (T_{\max}, M_{\max}, D_{\max})$ where:
\begin{itemize}[nosep]
  \item $T_{\max} \in \mathbb{N}$ is the maximum number of execution steps,
  \item $M_{\max} \in \mathbb{N}$ is the maximum memory consumption in bytes,
  \item $D_{\max} \in \mathbb{N}$ is the maximum call stack depth.
\end{itemize}
\end{definition}

\begin{theorem}[Computability of Resource Certificates]
\label{thm:mpr}
For every well-formed constrained-form program $P_C$, the resource certificate
$\mathcal{R}(P_C)$ is computable in time polynomial in the size of $P_C$.
\end{theorem}

\begin{proof}[Proof sketch]
We proceed by structural induction on the program.

\textbf{Base case:} For a program consisting of a single statement $s$ with
no loops or calls, $T_{\max}(s) = c_s$ (a constant depending on the
statement kind), $M_{\max}(s)$ is the size of variables bound by $s$, and
$D_{\max}(s) = 0$.

\textbf{Sequential composition:} For $s_1; s_2$,
\begin{align*}
  T_{\max}(s_1; s_2) &= T_{\max}(s_1) + T_{\max}(s_2), \\
  M_{\max}(s_1; s_2) &= \max(M_{\max}(s_1), M_{\max}(s_2)), \\
  D_{\max}(s_1; s_2) &= \max(D_{\max}(s_1), D_{\max}(s_2)).
\end{align*}

\textbf{Conditional:} For $\texttt{if}\;e\;\{s_t\}\;\texttt{else}\;\{s_f\}$,
\begin{align*}
  T_{\max} &= c_{\text{cond}} + \max(T_{\max}(s_t), T_{\max}(s_f)), \\
  M_{\max} &= \max(M_{\max}(s_t), M_{\max}(s_f)), \\
  D_{\max} &= \max(D_{\max}(s_t), D_{\max}(s_f)).
\end{align*}

\textbf{Bounded loop:} For $\texttt{for}\;i\;\texttt{in}\;0\texttt{..}N\;\{s_b\}$
where $N$ is a compile-time constant,
\begin{align*}
  T_{\max} &= c_{\text{init}} + N \cdot (T_{\max}(s_b) + c_{\text{iter}}), \\
  M_{\max} &= M_{\text{iter}} + M_{\max}(s_b), \\
  D_{\max} &= D_{\max}(s_b).
\end{align*}

\textbf{Function call:}  Because the call graph is a DAG (verified
statically), we can topologically sort functions and compute resource
certificates bottom-up.  A call to function $f$ with certificate
$\mathcal{R}(f)$ contributes $T_{\max}(f)$ to time, $M_{\max}(f)$ to
memory (added to the caller's frame), and $D_{\max}(f) + 1$ to depth.

The topological sort ensures that every function's callees have already been
analysed.  The absence of cycles makes this well-founded.  Each step of the
analysis is constant-time in the size of the statement, and the topological
sort is $\bigO(|V| + |E|)$ in the call graph.  The overall analysis is
therefore polynomial.
\end{proof}

\subsection{Static Analysis Implementation}

The \obl{} compiler includes a static analyser that computes resource
certificates.  The analysis result for a program includes:

\begin{lstlisting}[language={},caption={Static analysis output for Fibonacci example}]
=== Resource Bounds (Static Guarantees) ===
Max loop iterations: 10
Max call depth: 0
Estimated memory: 24 bytes

=== Reversibility Analysis ===
All reversible operations properly balanced

=== Accountability Trace Coverage ===
Coverage: 40.0%
\end{lstlisting}

The analysis is implemented in the \texttt{static\_analyzer.ml} module, which
traverses the AST bottom-up, accumulating resource bounds at each node.  The
key data structure is:

\begin{lstlisting}[language=ocaml-ast,caption={Resource analysis result type (OCaml)}]
type analysis_result = {
  is_valid: bool;
  violations: constraint_violation list;
  max_loop_iterations: int64;
  max_call_depth: int;
  estimated_memory: int;  (* bytes *)
  reversibility_warnings: string list;
  trace_coverage: float;  (* 0.0 to 1.0 *)
}
\end{lstlisting}

\subsection{Comparison with WCET Analysis}

Worst-case execution time (WCET) analysis \citep{wilhelm2008wcet} for
conventional languages is an entire subfield of real-time systems research,
involving abstract interpretation of machine-level cache and pipeline
behaviour.  WCET analysis is necessary precisely because the source language
does not bound execution---the bounds must be inferred from the program and
the hardware.

In \obl{}, the situation is qualitatively different: the source language
\emph{syntactically constrains} execution to be bounded, and the bounds are
\emph{manifest} in the program text (as loop limits and the call-graph
structure).  Our static analysis computes these manifest bounds; it does not
need to infer them.  This eliminates the fundamental difficulty of WCET
analysis---the undecidability of loop bounds in Turing-complete
languages---and replaces it with a straightforward polynomial-time
computation.


% ============================================================================
% 5. REVERSIBILITY
% ============================================================================
\section{Reversibility by Construction}
\label{sec:reversibility}

\subsection{Reversibility Invariant}

\begin{definition}[Reversibility]
\label{def:reversibility}
An operation $\mathsf{op}$ with state transformer
$\mathsf{op} : (\Sigma \times \mathcal{T}) \to (\Sigma \times \mathcal{T})$
(where $\Sigma$ is the program state and $\mathcal{T}$ is the accountability
trace) is \emph{reversible} if there exists an inverse operation
$\inv{\mathsf{op}}$ such that for all $(\sigma, \tau)$ in the domain:
\[
  \inv{\mathsf{op}}(\mathsf{op}(\sigma, \tau)) = (\sigma, \tau')
\]
where $\tau' \supseteq \tau$ (the trace is extended, never truncated).
\end{definition}

The definition permits the trace to grow during reversal---reversal itself is
an accountable operation.  This is essential for audit compliance: rolling
back a transaction must itself be recorded.

\subsection{Operation Classification}

Constrained-form operations fall into three reversibility classes:

\begin{table}[htbp]
\centering
\caption{Reversibility classification of constrained-form operations}
\label{tab:reversibility}
\begin{tabular}{@{}lll@{}}
\toprule
\textbf{Class} & \textbf{Operations} & \textbf{Inverse} \\
\midrule
Self-inverse & \texttt{swap(a, b)}, \texttt{x \^{}= v}, \texttt{not(x)} & Same operation \\
Paired inverse & \texttt{incr(x, d)} / \texttt{decr(x, d)} & Partner operation \\
 & \texttt{push(s, v)} / \texttt{pop(s)} & Partner operation \\
Trace-assisted & \texttt{x = e} (assignment) & Restore from trace \\
\bottomrule
\end{tabular}
\end{table}

Self-inverse operations are their own inverses: applying them twice restores
the original state.  This is immediate for XOR-assign ($x \oplus v \oplus v
= x$) and swap.

Paired-inverse operations have distinct inverses: incrementing by $\delta$ is
reversed by decrementing by $\delta$.  The compiler verifies that paired
operations are well-formed (the inverse argument matches the forward
argument).

Trace-assisted operations are \emph{locally irreversible}---assignment
destroys the previous value---but become reversible through the
accountability trace, which records the overwritten value.  This is the key
mechanism that bridges classical imperative programming with reversibility.

\begin{theorem}[Constrained-Form Reversibility]
\label{thm:reversibility}
Every well-formed constrained-form program is reversible in the sense of
\cref{def:reversibility}.
\end{theorem}

\begin{proof}
By structural induction on statements.

\textbf{Self-inverse operations:}  $\inv{\texttt{swap}(a,b)} =
\texttt{swap}(a,b)$ and $\inv{(x \oplus v)} = (x \oplus v)$.  These satisfy
the identity law directly.

\textbf{Paired operations:}  $\inv{\texttt{incr}(x, \delta)} =
\texttt{decr}(x, \delta)$ and vice versa, with
$\texttt{decr}(\texttt{incr}(x, \delta), \delta) = x$ by integer arithmetic.

\textbf{Assignment:}  For $x := e$, the trace records $(\textit{old}(x),
\textit{new}(x))$.  The inverse reads $\textit{old}(x)$ from the trace and
restores it.

\textbf{Sequential composition:}  $\inv{(s_1; s_2)} = \inv{s_2}; \inv{s_1}$.
This reversal of order is standard in reversible computation.

\textbf{Conditional:}  The trace records which branch was taken.  The inverse
reads the branch indicator from the trace and executes the inverse of the
taken branch.

\textbf{Bounded loop:}  $\texttt{for}\;i\;\texttt{in}\;0\texttt{..}N\;\{s\}$
reverses to $\texttt{for}\;i\;\texttt{in}\;(N{-}1)\texttt{..}(-1)\;\{\inv{s}\}$.
The loop bounds are static constants, so the reversed loop also has static
bounds.

\textbf{Function call:}  Because the call graph is acyclic, induction on the
topological order of the call graph ensures that every called function has a
well-defined inverse.  The inverse of a call to $f$ is a call to $\inv{f}$.
\end{proof}

\subsection{The Reversible Debugger}

A direct consequence of program-level reversibility is the \obl{} reversible
debugger, which supports backward stepping through execution.  The debugger
maintains the accountability trace during forward execution and uses the
inverse semantics to step backward:

\begin{lstlisting}[language={},caption={Reversible debugger session}]
oblibeny> step       (* execute incr(c, 1) *)
  c: 0 -> 1
  trace: [checkpoint("start"), increment(c, 1)]

oblibeny> step       (* execute incr(c, 1) *)
  c: 1 -> 2
  trace: [..., increment(c, 1)]

oblibeny> back       (* REVERSE: decr(c, 1) *)
  c: 2 -> 1
  trace: [..., reverse(increment(c, 1))]

oblibeny> back       (* REVERSE: decr(c, 1) *)
  c: 1 -> 0
  trace: [..., reverse(increment(c, 1))]
\end{lstlisting}

The debugger's backward step is not a heuristic or a ``replay from
checkpoint'' mechanism---it is a direct execution of the inverse operation.
This is only possible because (a)~every operation has a defined inverse and
(b)~the trace provides sufficient information to reconstruct destroyed state.

\subsection{Transaction Rollback Without Undo Logic}

In conventional systems, implementing transaction rollback requires
application-specific undo logic: if a transaction writes to a database, the
rollback must issue compensating writes.  This undo logic is a separate code
path that must be written, tested, and maintained in parallel with the forward
logic.

In \obl{}, rollback is automatic.  Given a checkpoint~$c$ in the trace, the
runtime can reverse all operations since~$c$ by executing their inverses in
reverse order.  No application-specific undo logic is needed.  This
eliminates an entire category of bugs---incorrect or incomplete rollback---by
construction.


% ============================================================================
% 6. THE ACCOUNTABILITY TRACE
% ============================================================================
\section{The Accountability Trace}
\label{sec:trace}

Every execution of a constrained-form program produces an
\emph{accountability trace}: an append-only, cryptographically-chained log of
all operations performed.

\subsection{Trace Structure}

Each trace entry is a record:
\[
  e_i = (\mathit{seq}_i, \; \mathit{op}_i, \; \mathit{inputs}_i, \;
         \mathit{outputs}_i, \; \mathit{actor}_i, \; h_{i-1})
\]
where $\mathit{seq}_i$ is a monotonically-increasing sequence number,
$\mathit{op}_i$ identifies the operation, $\mathit{inputs}_i$ and
$\mathit{outputs}_i$ record the operation's arguments and results (hashed or
literal, depending on configuration), $\mathit{actor}_i$ identifies the
executing principal, and $h_{i-1} = H(e_{i-1})$ is the cryptographic hash of
the previous entry.

The hash chain provides tamper evidence: modifying any entry invalidates all
subsequent hashes.  The trace is sufficient to:
\begin{enumerate}[nosep]
  \item \textbf{Replay} the computation: execute each $\mathit{op}_i$ with
    $\mathit{inputs}_i$ and verify that $\mathit{outputs}_i$ matches.
  \item \textbf{Reverse} the computation: execute $\inv{\mathit{op}_i}$
    for each entry in reverse order.
  \item \textbf{Audit} the computation: verify the hash chain integrity and
    check that all operations conform to policy.
\end{enumerate}

\subsection{Trace Completeness}

\begin{theorem}[Trace Sufficiency]
For every constrained-form execution, the accountability trace is sufficient
to reconstruct any prior program state.
\end{theorem}

\begin{proof}
By \cref{thm:reversibility}, every operation is reversible given the trace.
By the sequential reversal property, the entire execution can be reversed to
any prior state.  The trace records all information needed for reversal
(including overwritten values for trace-assisted operations).  Therefore, the
trace, combined with the final state, suffices to reconstruct any
intermediate state.
\end{proof}

\subsection{Bounded Trace Size}

Because the program terminates in at most $T_{\max}$ steps (by
\cref{thm:mpr}), the trace contains at most $T_{\max}$ entries.  If each
entry has bounded size~$c_e$ (determined by the maximum size of operation
arguments), the total trace size is bounded by $T_{\max} \cdot c_e$.  This
bound is \emph{statically computable}, meaning that the trace's storage
requirements can be allocated in advance---a critical property for HSM
firmware where dynamic allocation may be prohibited.

\subsection{Cryptographic Binding}

The hash chain uses a configurable hash function.  For post-quantum
resistance, the default is SHA-3-256; for classical deployments, BLAKE3 is
supported.  The trace can optionally be signed by the executing principal
using ML-DSA (Dilithium) for non-repudiation.

The trace integrates with \obl{}'s external accountability protocol (MAA,
Mutually-Assured Accountability): at checkpoints, the current trace hash
can be committed to an external ledger, providing third-party verifiable
evidence of execution history.


% ============================================================================
% 7. IMPLEMENTATION
% ============================================================================
\section{The OCaml Compiler}
\label{sec:implementation}

\subsection{Architecture}

The \obl{} compiler is implemented in approximately 5,200 lines of OCaml,
using the Menhir parser generator for the constrained-form grammar and a
hand-written recursive-descent parser for the factory form.  The
implementation comprises:

\begin{table}[htbp]
\centering
\caption{Compiler module structure}
\label{tab:modules}
\begin{tabular}{@{}llr@{}}
\toprule
\textbf{Module} & \textbf{Purpose} & \textbf{LOC} \\
\midrule
\texttt{lexer.ml} & Tokenisation (sedlex) & 180 \\
\texttt{parser.mly} & Menhir grammar & 220 \\
\texttt{ast.ml} & Abstract syntax tree & 170 \\
\texttt{typecheck.ml} & Type inference and checking & 420 \\
\texttt{constrained\_check.ml} & Turing-incompleteness validation & 200 \\
\texttt{static\_analyzer.ml} & Resource bounds computation & 350 \\
\texttt{eval.ml} & Constrained-form evaluator & 480 \\
\texttt{trace.ml} & Accountability trace runtime & 180 \\
\texttt{debugger.ml} & Reversible debugger & 300 \\
\texttt{profiler.ml} & Resource profiling & 250 \\
\texttt{lsp\_*.ml} & Language Server Protocol & 789 \\
\texttt{errors.ml}, \texttt{location.ml} & Diagnostics & 150 \\
\bottomrule
\end{tabular}
\end{table}

\subsection{The Constrained-Form Checker}

The constrained-form checker (\texttt{constrained\_check.ml}) is the
component that enforces Turing-incompleteness.  It operates in two phases:

\textbf{Phase 1: Local checks.}  For each function declaration, the checker
verifies:
\begin{itemize}[nosep]
  \item No direct self-recursion: the function body does not contain a call
    to the function's own name.
  \item All \texttt{for}-range bounds are compile-time constants.
  \item No forbidden keywords (\texttt{while}, \texttt{loop}) appear in the
    token stream---these are rejected by the lexer.
\end{itemize}

\textbf{Phase 2: Global call-graph check.}  The checker builds the call graph
$G = (V, E)$ where $V$ is the set of function names and $(f, g) \in E$ if
$f$'s body contains a call to $g$.  It then checks that $G$ is acyclic via
depth-first search, detecting strongly-connected components.

If any violation is detected, the compiler reports a diagnostic with the
violation location and a clear error message:

\begin{lstlisting}[language={},caption={Constrained-form violation diagnostics}]
src/main.obl:12:5: error: recursive call from 'fibonacci'
  to 'fibonacci' not allowed in constrained form

src/main.obl:20:5: error: for-range bound 'n' must be a
  static constant

error: cyclic call graph detected: ping -> pong -> ping
\end{lstlisting}

\subsection{AST Design}

The AST (\texttt{ast.ml}) is designed to represent both forms.  The
constrained-form restriction is that only a subset of the AST constructors is
valid: the \texttt{SForRange} constructor requires \texttt{int64} constants
for its bounds (not arbitrary expressions), and the constraint violation types
enumerate the precise ways a program can fail to be constrained-form-valid.

The key AST types are:

\begin{lstlisting}[language=ocaml-ast,caption={Core AST types (excerpt)}]
type stmt_desc =
  | SLet of string * typ option * expr
  | SLetMut of string * typ option * expr
  | SAssign of string * expr
  | SIf of expr * stmt list * stmt list
  | SForRange of string * int64 * int64 * stmt list
  | SSwap of string * string
  | SIncr of string * expr
  | SDecr of string * expr
  | SXorAssign of string * expr
  | STrace of string * expr list
  | SCheckpoint of string
  | SAssertInvariant of expr * string
\end{lstlisting}

Note the \texttt{SForRange} constructor takes \texttt{int64} bounds directly,
not \texttt{expr}.  This is an example of \emph{encoding the restriction in
the type}: it is impossible to construct a \texttt{SForRange} node with
dynamic bounds, because the constructor does not accept expressions for
bounds.

\subsection{Idris2 ABI Proofs}

The foreign function interface between \obl{} and the post-quantum
cryptographic libraries uses an Idris2-defined ABI:

\begin{lstlisting}[language={},caption={Idris2 ABI definition (excerpt from Crypto.idr)}]
-- Proven properties of the cryptographic ABI:
-- 1. Key sizes are correct for each algorithm
-- 2. Buffer sizes match algorithm requirements
-- 3. Return codes are exhaustively handled
-- 4. No undefined behaviour in any calling convention
\end{lstlisting}

The Idris2 types encode invariants that are verified at compile time using
dependent types.  For example, the type of a signing function encodes that the
signature buffer must have exactly the size specified by the algorithm
parameter---a size mismatch is a type error, not a runtime buffer overflow.

\subsection{Zig FFI Layer}

The Zig FFI provides the actual implementation of cryptographic operations,
binding to liboqs (for ML-KEM, ML-DSA, SLH-DSA) and libsodium (for classical
cryptography and utilities).  Zig was chosen for this layer because it:

\begin{itemize}[nosep]
  \item Produces C-ABI-compatible binaries with no runtime dependencies.
  \item Provides memory safety by default (no hidden allocations, explicit
    error handling).
  \item Supports cross-compilation to HSM target architectures (ARM
    Cortex-M, RISC-V) from the build system.
  \item Has zero-cost abstractions---the FFI layer adds no overhead over a
    direct C binding.
\end{itemize}

The architecture is:

\begin{center}
\begin{tikzpicture}[
  box/.style={draw, rounded corners=3pt, minimum width=3.2cm,
              minimum height=0.8cm, align=center, font=\small},
  arr/.style={-{Stealth[length=4pt]}, thick}
]
  \node[box, fill=yellow!15] (idris) {Idris2 ABI\\(formal proofs)};
  \node[box, fill=cyan!10, below=0.8cm of idris] (cheader) {Generated C Headers};
  \node[box, fill=green!10, below=0.8cm of cheader] (zig) {Zig FFI\\(implementation)};
  \node[box, fill=red!8, below left=0.8cm and -0.2cm of zig] (liboqs) {liboqs\\(PQ crypto)};
  \node[box, fill=red!8, below right=0.8cm and -0.2cm of zig] (sodium) {libsodium\\(classical)};

  \draw[arr] (idris) -- node[right, font=\footnotesize\itshape] {generates} (cheader);
  \draw[arr] (cheader) -- node[right, font=\footnotesize\itshape] {implements} (zig);
  \draw[arr] (zig) -- (liboqs);
  \draw[arr] (zig) -- (sodium);
\end{tikzpicture}
\end{center}


% ============================================================================
% 8. CASE STUDIES
% ============================================================================
\section{Case Studies}
\label{sec:case-studies}

\subsection{HSM Key Management Firmware}

We evaluate \obl{} against the requirements of a hypothetical HSM key
management module---a firmware component responsible for generating,
storing, rotating, and destroying cryptographic keys within a FIPS~140-3
Level~3 boundary.

\subsubsection{Requirements}

The HSM firmware must satisfy:
\begin{enumerate}[nosep,label=(H\arabic*)]
  \item \textbf{Termination:}  Every operation must complete within a bounded
    time, specified in the HSM's security policy.  An operation that exceeds
    its time budget triggers a hardware reset.
  \item \textbf{Auditability:}  Every key lifecycle event (generation,
    export, rotation, destruction) must be recorded in a tamper-evident log.
  \item \textbf{Reversibility:}  Key rotation must support rollback in case
    of partial failure (e.g., if the new key is generated but the old key has
    not yet been marked as deprecated).
  \item \textbf{Resource bounds:}  The firmware must not exceed its allocated
    SRAM (typically 256\,KB--1\,MB), and must not perform dynamic allocation
    after initialisation.
  \item \textbf{Post-quantum readiness:}  The module must support ML-KEM and
    ML-DSA in addition to classical algorithms.
\end{enumerate}

\subsubsection{\obl{} Solution}

\begin{lstlisting}[language=oblibeny,caption={HSM key rotation in constrained form}]
fn rotate_key(
  slot: u32, new_key: &KeyMaterial
) -> RotationResult {
  checkpoint("pre_rotation");

  // Record old key hash for audit
  let old_hash: Hash = key_hash(slot);
  trace("rotation_start", slot, old_hash);

  // Reversible state transition
  let mut backup: KeyMaterial = read_key(slot);
  swap(backup, new_key);  // Self-inverse
  write_key(slot, new_key);
  trace("key_written", slot, key_hash(slot));

  // Mark old key for zeroisation (bounded loop)
  for i in 0..32 {
    incr(backup.bytes[i], 0);  // Touch for trace
  }

  checkpoint("post_rotation");
  assert_invariant(
    key_hash(slot) != old_hash,
    "New key must differ from old"
  );

  RotationResult::Ok
}
\end{lstlisting}

The static analyser produces a resource certificate for this function:
$T_{\max} = 847$ steps, $M_{\max} = 312$ bytes, $D_{\max} = 1$.  These
bounds are checked against the HSM's security policy \emph{at compile time}.
If the bounds exceed the policy limits, the firmware fails to compile.

If the rotation fails partway (e.g., \texttt{write\_key} returns an error),
the runtime reverses all operations since the ``pre\_rotation'' checkpoint,
restoring the old key without any application-specific rollback logic.

\subsection{Secure Enclave Attestation}

We consider a TrustZone enclave that performs remote attestation---proving to
a verifier that the enclave is running expected code on genuine hardware.

\subsubsection{Challenge}

The attestation protocol involves:
\begin{enumerate}[nosep]
  \item Receiving a nonce from the verifier.
  \item Measuring the enclave's code and data.
  \item Signing the measurement with a hardware-bound key.
  \item Returning the signed attestation to the verifier.
\end{enumerate}

Each step must terminate; the attestation must be auditable; and the signing
must use post-quantum algorithms for long-term security.

\subsubsection{\obl{} Solution}

\begin{lstlisting}[language=oblibeny,caption={Enclave attestation in constrained form}]
fn attest(nonce: &[u8; 32]) -> Attestation {
  checkpoint("attestation_start");

  // Measure enclave state (bounded by region count)
  let mut measurement: Hash = initial_hash();
  for region in 0..8 {
    let region_hash: Hash = hash_region(region);
    measurement = combine_hashes(
      measurement, region_hash
    );
    trace("measured_region", region, region_hash);
  }

  // Incorporate nonce (prevents replay)
  measurement = combine_hashes(measurement, nonce);
  trace("nonce_incorporated", nonce);

  // Sign with ML-DSA (post-quantum)
  let sig: Signature = ml_dsa_sign(
    hardware_key(), measurement
  );
  trace("attestation_signed", sig);

  checkpoint("attestation_complete");

  Attestation { measurement, sig, nonce: *nonce }
}
\end{lstlisting}

The static bound on this function is determined entirely by the constant~8
(the number of memory regions to measure) and the constant costs of the
cryptographic operations (known from the FFI specifications).  The
accountability trace provides a complete record of what was measured, in what
order, and with what result---valuable for both debugging and compliance.

\subsection{Quantitative Evaluation}

\begin{table}[htbp]
\centering
\caption{Comparison of guarantee mechanisms across approaches}
\label{tab:comparison}
\begin{tabular}{@{}lccc@{}}
\toprule
\textbf{Property} & \textbf{C + WCET} & \textbf{Agda/Idris2} & \textbf{\obl{}} \\
\midrule
Termination guaranteed & Per-program & Per-program & All programs \\
Resource bounds static & Approx. & Manual & Automatic \\
Reversibility & Manual undo & Not built-in & By construction \\
Accountability trace & Manual logging & Not built-in & By construction \\
New program effort & Re-verify & Re-check & Zero \\
Post-quantum FFI & Manual binding & Possible & Verified \\
\bottomrule
\end{tabular}
\end{table}

The key advantage of \obl{}'s approach is the ``new program effort'' row: when
a developer writes a new constrained-form program, it is \emph{automatically}
terminating, resource-bounded, reversible, and auditable.  No per-program
verification work is needed.  In contrast, both the ``C + WCET'' approach and
the dependently-typed approach require effort proportional to the program's
complexity.


% ============================================================================
% 9. DISCUSSION
% ============================================================================
\section{Discussion}
\label{sec:discussion}

\subsection{When Turing-Incompleteness is a Feature}

The conventional wisdom in programming language design treats
Turing-completeness as a baseline requirement.  A language that is not
Turing-complete is considered ``limited'' or ``toy.''  We argue that this
perspective is inverted for the deployment contexts we consider.

In an HSM, the \emph{ability} to write a non-terminating program is not a
feature---it is a hazard.  Every non-terminating program that \emph{could} be
written is a program that a developer \emph{might} accidentally write, and
that a review process \emph{might} fail to catch.  Eliminating the hazard at
the language level is strictly superior to mitigating it through process.

This parallels the evolution of memory safety.  C permits arbitrary pointer
arithmetic; Rust prevents it through the type system.  The Rust approach is
not ``limiting''---it is a deliberate engineering trade-off that eliminates an
entire class of bugs.  \obl{} makes an analogous trade-off for termination,
reversibility, and accountability.

\subsection{Expressiveness of the Constrained Form}

A natural question is: what useful programs \emph{cannot} be written in
constrained form?

Any program that requires unbounded iteration over input data cannot be
directly expressed.  This includes, for example, a general-purpose string
search over an input of unknown length.  However, in HSM and enclave
contexts, input sizes are typically bounded by hardware constraints (key sizes
are fixed, message buffers have maximum lengths, and the number of key slots
is finite).  The factory form can generate constrained-form programs
parameterised by these hardware-specific bounds.

Programs that are inherently Turing-complete---such as interpreters for
Turing-complete languages---cannot be expressed in constrained form.  This is
by design: deploying an interpreter inside an HSM would defeat the
termination guarantee.

\subsection{The Factory as a Safety Contract}

The dual-form model inverts the traditional safety analysis workflow.
Instead of:
\begin{quote}
  \emph{Write program $\to$ verify program $\to$ deploy}
\end{quote}
the workflow becomes:
\begin{quote}
  \emph{Write factory $\to$ factory produces constrained program
  $\to$ constrained program is safe by construction $\to$ deploy}
\end{quote}

The verification burden shifts from the deployed artifact to the factory.
And the factory does not need to be verified for termination or
reversibility---it runs in a development environment where these properties
are not required.  This is a form of \emph{responsibility separation}:
the factory is responsible for correctness (producing the right program), and
the constrained form is responsible for safety (executing within bounds).

\subsection{Relationship to Smart Contract Languages}

Several smart contract languages have adopted Turing-incompleteness or
bounded computation for related reasons.  The Bitcoin Script language is
intentionally non-Turing-complete, and Ethereum's gas mechanism provides
\emph{economic} bounds on computation (though not static bounds).
Move~\citep{blackshear2019move}, developed for the Diem blockchain, uses a
restricted form of references and a bytecode verifier to ensure safety
properties.

\obl{} differs from these in two respects.  First, \obl{} targets
\emph{hardware} security rather than blockchain execution---the constraint
motivation is physical safety, not economic cost.  Second, \obl{}'s dual-form
model provides a Turing-complete factory, whereas smart contract languages
typically have a single restricted form.

\subsection{Limitations}

\begin{enumerate}[nosep]
  \item \textbf{Factory termination:}  We do not guarantee that the factory
    form terminates.  A non-terminating factory program causes the compiler
    to time out, which is acceptable in a development environment but means
    that the ``time to compile'' is not statically bounded.
  \item \textbf{FFI safety:}  The Zig FFI layer is outside the
    constrained-form guarantee.  A bug in the FFI implementation could
    introduce non-termination or irreversibility.  The Idris2 proofs mitigate
    this at the type level, but do not cover implementation errors.
  \item \textbf{Expressiveness ceiling:}  Some algorithms that are naturally
    expressed with unbounded recursion (e.g., parsing context-free grammars)
    require reformulation as bounded iteration with explicit stack management.
    This is possible but increases code complexity.
  \item \textbf{Trace overhead:}  Recording every operation in the
    accountability trace imposes runtime overhead.  For most HSM operations
    (which are dominated by cryptographic computation), this overhead is
    negligible, but for tight inner loops it may be significant.
\end{enumerate}


% ============================================================================
% 10. RELATED WORK
% ============================================================================
\section{Related Work}
\label{sec:related}

\subsection{Total Languages}

Turner's Total Functional Programming \citep{turner2004total} is the
philosophical ancestor of \obl{}'s constrained form.  Turner argued that
totality is not merely an academic concern but a practical one: total programs
are easier to reason about, test, and compose.  \obl{} extends Turner's
argument to imperative programming in safety-critical contexts.

The Charity language \citep{cockett1992charity} enforced totality through
categorical semantics.  Every Charity program corresponds to a morphism in a
distributive category, ensuring totality by construction.  \obl{} achieves a
similar effect through syntactic restriction rather than categorical
foundations.

System~F$_\omega$ and the Calculus of Constructions are Turing-incomplete
type theories that serve as the core of dependently-typed proof assistants
(Coq, Agda, Lean).  These systems achieve termination through structural
recursion on inductive types.  \obl{}'s constrained form is less expressive
(no recursion at all) but simpler to implement and verify.

\subsection{Reversible Languages}

Janus \citep{yokoyama2007reversible} is the most developed reversible
imperative language.  Janus restricts assignment to reversible operations and
uses assertion-guarded conditionals to ensure backward determinism.  R-WHILE
\citep{gluck2012rwhile} and R-CORE extend reversibility to structured
iteration.

\obl{} differs from these languages in combining reversibility with
Turing-incompleteness: backward execution in \obl{} is also guaranteed to
terminate, whereas backward execution of a Janus program may diverge if the
forward execution diverged.

The Theseus language \citep{james2014theseus} explored reversibility in a
functional setting using type isomorphisms.  \obl{}'s trace-assisted
reversibility is more pragmatic but less elegant.

\subsection{Multi-Stage Programming}

MetaOCaml \citep{kiselyov2014metaocaml} and Template Haskell
\citep{sheard2002template} are the best-known multi-stage programming
systems.  Both allow compile-time code generation within a
Turing-complete language.  \obl{}'s dual-form model is a multi-stage system
where the stages have \emph{different computational power}---to our knowledge,
this asymmetric staging is novel.

Terra \citep{devito2013terra} embeds a low-level language (for GPU and
systems programming) within Lua (for metaprogramming).  The relationship is
analogous to \obl{}'s factory/constrained split, though Terra does not enforce
termination or reversibility in the low-level language.

\subsection{Domain-Specific Languages for Security}

The Cryptol language \citep{lewis2003cryptol} is a domain-specific language
for cryptographic algorithm specification.  Cryptol enforces finite-bit-width
types and bounded streams, providing some of the same termination guarantees
as \obl{}.  However, Cryptol lacks reversibility and accountability.

SAW (Software Analysis Workbench) \citep{dockins2016saw} provides formal
verification of C and LLVM code against Cryptol specifications.  This is a
post-hoc verification approach, complementary to \obl{}'s by-construction
approach.

F* \citep{swamy2016dependent} is a dependently-typed language for verified
security-critical programming.  F*'s effect system can express resource
bounds, but these are verified per-program rather than guaranteed by language
restriction.

\subsection{Resource-Bounded Computation}

Implicit computational complexity (ICC) \citep{dal2010linear} studies
characterisations of complexity classes through type systems.  Light linear
logic \citep{girard1998light} and related systems characterise polynomial
time through linear type restrictions.  \obl{}'s resource bounds are simpler
than ICC characterisations (bounded loops and acyclic call graphs rather than
linear type discipline) but are sufficient for the target domain.

The Celf logical framework \citep{schack2015celf} uses linear logic for
resource-aware specification.  \obl{}'s approach is more pragmatic: resources
are bounded by the syntactic structure (loop bounds and call-graph depth)
rather than by a logical type discipline.


% ============================================================================
% 11. CONCLUSION
% ============================================================================
\section{Conclusion}
\label{sec:conclusion}

We have presented \obl{}, a dual-form programming language that separates
Turing-complete metaprogramming (at compile time) from Turing-incomplete
constrained execution (at runtime).  The constrained form guarantees
termination, bounded resource consumption, full reversibility, and complete
accountability---not through per-program verification, but through language
design that makes violations \emph{inexpressible}.

The dual-form model recovers developer expressiveness without sacrificing
deployment safety: the factory form can compute arbitrary functions to produce
the constrained artifact, and the constrained artifact is what faces the
hostile environment.  This separation of concerns---correctness in the
factory, safety in the constrained form---is a principled engineering
trade-off suited to a class of deployment targets where the cost of
non-termination is measured not in user frustration but in hardware
destruction, regulatory penalty, and physical safety.

The compiler is implemented and production-ready, with formal ABI proofs in
Idris2, post-quantum cryptographic support via Zig FFI, and a complete
toolchain including LSP server, reversible debugger, static analyser, and
container deployment.  The full implementation, specification, and example
programs are available at
\url{https://github.com/hyperpolymath/oblibeny}.

\paragraph{Future work.}
We plan to extend the constrained form with \emph{dependent types} that
encode resource bounds in the type system (e.g., a function type
$\texttt{fn}(n : \texttt{u32}) \to_{\leq 10n + 5} \texttt{u32}$ that
declares its resource consumption as a function of its input), and to develop
a RISC-V backend targeting the open-source HSM reference platform.  We also
plan to formalise the full dual-form model in Agda, proving the metatheoretic
properties (soundness of the resource analysis, completeness of the
reversibility semantics) mechanically.


% ============================================================================
% ACKNOWLEDGMENTS
% ============================================================================
\section*{Acknowledgements}

The author thanks the developers of OCaml, Idris2, Zig, liboqs, and
libsodium, whose tools form the implementation substrate of this work.  The
design of \obl{}'s reversibility primitives was influenced by conversations
about Janus and reversible computing at the intersection of programming
languages and thermodynamic theory.


% ============================================================================
% REFERENCES
% ============================================================================

\begin{thebibliography}{99}

\bibitem[Bennett(1989)]{bennett1989timespace}
Bennett, C.~H. (1989).
\newblock Time/space trade-offs for reversible computation.
\newblock \emph{SIAM Journal on Computing}, 18(4):766--776.

\bibitem[Blackshear et~al.(2019)]{blackshear2019move}
Blackshear, S., Cheng, E., Dill, D.~L., Gao, V., Maurer, B., Nowacki, T., Pott, A., Qadeer, S., Rain, Russinoff, D., Stephlin, S., Zakian, T., and Zhou, R. (2019).
\newblock Move: A language with programmable resources.
\newblock \emph{Libra Technical Whitepaper}.

\bibitem[Brady(2021)]{brady2021idris2}
Brady, E. (2021).
\newblock \emph{Idris 2: Quantitative Type Theory in Practice}.
\newblock In \emph{ECOOP 2021}, LIPIcs 194.

\bibitem[Cockett and Fukushima(1992)]{cockett1992charity}
Cockett, R. and Fukushima, T. (1992).
\newblock About Charity.
\newblock Technical Report 92/480/18, University of Calgary.

\bibitem[The Coq Development Team(2024)]{coq2024}
The Coq Development Team (2024).
\newblock \emph{The Coq Proof Assistant Reference Manual}.
\newblock Version 8.19.

\bibitem[Costan and Devadas(2016)]{costan2016intel}
Costan, V. and Devadas, S. (2016).
\newblock Intel SGX explained.
\newblock \emph{Cryptology ePrint Archive}, Report 2016/086.

\bibitem[Dal~Lago(2010)]{dal2010linear}
Dal~Lago, U. (2010).
\newblock A short introduction to implicit computational complexity.
\newblock In \emph{Lectures on Logic and Computation}, pp.~89--109. Springer.

\bibitem[DeVito et~al.(2013)]{devito2013terra}
DeVito, Z., Hegarty, J., Aiken, A., Hanrahan, P., and Vitek, J. (2013).
\newblock Terra: A multi-stage language for high-performance computing.
\newblock In \emph{PLDI 2013}, pp.~105--116.

\bibitem[Dockins et~al.(2016)]{dockins2016saw}
Dockins, R., Foltzer, A., Hendrix, J., Huffman, B., McNamee, D., and Tomb, A. (2016).
\newblock Constructing semantic models of programs with the Software Analysis Workbench.
\newblock In \emph{VSTTE 2016}, pp.~56--72.

\bibitem[{FIPS 140-3}(2019)]{fips1403}
{National Institute of Standards and Technology} (2019).
\newblock Security Requirements for Cryptographic Modules.
\newblock \emph{Federal Information Processing Standard} 140-3.

\bibitem[Girard(1998)]{girard1998light}
Girard, J.-Y. (1998).
\newblock Light linear logic.
\newblock \emph{Information and Computation}, 143(2):175--204.

\bibitem[Gl\"{u}ck and Yokoyama(2012)]{gluck2012rwhile}
Gl\"{u}ck, R. and Yokoyama, T. (2012).
\newblock A reversible programming language and its invertible self-interpreter.
\newblock In \emph{PEPM 2012}, pp.~144--153.

\bibitem[James and Sabry(2014)]{james2014theseus}
James, R.~P. and Sabry, A. (2014).
\newblock Theseus: A high-level approach to reversible computing.
\newblock \emph{Journal of Functional Programming}, 24(2--3):341--401.

\bibitem[Kiselyov(2014)]{kiselyov2014metaocaml}
Kiselyov, O. (2014).
\newblock The design and implementation of BER MetaOCaml.
\newblock In \emph{FLOPS 2014}, pp.~86--102.

\bibitem[Landauer(1961)]{landauer1961irreversibility}
Landauer, R. (1961).
\newblock Irreversibility and heat generation in the computing process.
\newblock \emph{IBM Journal of Research and Development}, 5(3):183--191.

\bibitem[Lewis et~al.(2003)]{lewis2003cryptol}
Lewis, J.~R., Martin, B., Goerigk, W., and Strake, M. (2003).
\newblock Cryptol: High assurance, retargetable crypto development and validation.
\newblock In \emph{MILCOM 2003}, pp.~820--825.

\bibitem[{Common Criteria}(2022)]{commoncriteria}
{Common Criteria Recognition Arrangement} (2022).
\newblock \emph{Common Criteria for Information Technology Security Evaluation}.
\newblock Version 3.1, Revision 5.

\bibitem[{NIST}(2024)]{nist2024pqc}
{National Institute of Standards and Technology} (2024).
\newblock Post-Quantum Cryptography Standardization.
\newblock FIPS 203 (ML-KEM), FIPS 204 (ML-DSA), FIPS 205 (SLH-DSA).

\bibitem[Norell(2007)]{norell2007towards}
Norell, U. (2007).
\newblock \emph{Towards a Practical Programming Language Based on Dependent Type Theory}.
\newblock PhD thesis, Chalmers University of Technology.

\bibitem[Pinto and Santos(2019)]{pinto2019demystifying}
Pinto, S. and Santos, N. (2019).
\newblock Demystifying Arm TrustZone: A comprehensive survey.
\newblock \emph{ACM Computing Surveys}, 51(6):1--36.

\bibitem[Schack-Nielsen and Sch\"{u}rmann(2015)]{schack2015celf}
Schack-Nielsen, A. and Sch\"{u}rmann, C. (2015).
\newblock Celf -- a logical framework for deductive and concurrent systems.
\newblock \emph{Journal of Automated Reasoning}, 55(1):1--32.

\bibitem[Sheard and Jones(2002)]{sheard2002template}
Sheard, T. and Jones, S.~P. (2002).
\newblock Template meta-programming for Haskell.
\newblock In \emph{Haskell Workshop 2002}, pp.~1--16.

\bibitem[Stebila and Mosca(2024)]{stebila2024liboqs}
Stebila, D. and Mosca, M. (2024).
\newblock Open Quantum Safe: Software for prototyping quantum-resistant cryptography.
\newblock \url{https://openquantumsafe.org/}.

\bibitem[Swamy et~al.(2016)]{swamy2016dependent}
Swamy, N., Hri\c{t}cu, C., Keller, C., Rastogi, A., Delignat-Lavaud, A., Forest, S., Bhargavan, K., Fournet, C., Strub, P.-Y., Kohlweiss, M., Zinzindohoue, J.-K., and Zanella-B\'{e}guelin, S. (2016).
\newblock Dependent types and multi-monadic effects in F*.
\newblock In \emph{POPL 2016}, pp.~256--270.

\bibitem[Taha(2004)]{taha2004gentle}
Taha, W. (2004).
\newblock A gentle introduction to multi-stage programming.
\newblock In \emph{Domain-Specific Program Generation}, LNCS 3016, pp.~30--50.

\bibitem[Turner(2004)]{turner2004total}
Turner, D.~A. (2004).
\newblock Total functional programming.
\newblock \emph{Journal of Universal Computer Science}, 10(7):751--768.

\bibitem[Wilhelm et~al.(2008)]{wilhelm2008wcet}
Wilhelm, R., Engblom, J., Ermedahl, A., Holsti, N., Thesing, S., Whalley, D., Bernat, G., Ferdinand, C., Heckmann, R., Mitra, T., Mueller, F., Puaut, I., Puschner, P., Staschulat, J., and Stenstr\"{o}m, P. (2008).
\newblock The worst-case execution-time problem---overview of methods and survey of tools.
\newblock \emph{ACM Transactions on Embedded Computing Systems}, 7(3):1--53.

\bibitem[Yokoyama et~al.(2007)]{yokoyama2007reversible}
Yokoyama, T., Axelsen, H.~B., and Gl\"{u}ck, R. (2007).
\newblock Principles of a reversible programming language.
\newblock In \emph{CF 2007}, pp.~43--54.

\end{thebibliography}


% ============================================================================
% APPENDIX
% ============================================================================
\appendix

\section{Full Constrained-Form Grammar}
\label{app:grammar}

The following is the complete context-free grammar for the constrained form,
in EBNF notation.  Terminal symbols are in \texttt{monospace}; nonterminals
are in \textit{italics}.

\begin{lstlisting}[language={},basicstyle=\ttfamily\footnotesize,frame=single,
  backgroundcolor=\color{gray!3}]
program     = { declaration } ;

declaration = function_decl
            | struct_decl
            | const_decl ;

function_decl = "fn" IDENT "(" [ param_list ] ")"
                "->" type "{" { statement } "}" ;

struct_decl = "struct" IDENT "{"
              { IDENT ":" type "," } "}" ;

const_decl  = "const" IDENT ":" type "=" expr ";" ;

param_list  = param { "," param } ;
param       = IDENT ":" type ;

type        = "i32" | "i64" | "u32" | "u64"
            | "bool" | "()"
            | "[" type ";" INT_LIT "]"
            | "&" type
            | IDENT ;

statement   = let_stmt | let_mut_stmt | assign_stmt
            | if_stmt | for_range_stmt
            | swap_stmt | incr_stmt | decr_stmt
            | xor_assign_stmt
            | trace_stmt | checkpoint_stmt
            | assert_stmt
            | expr_stmt | return_stmt ;

let_stmt       = "let" IDENT ":" type "=" expr ";" ;
let_mut_stmt   = "let" "mut" IDENT ":" type "=" expr ";" ;
assign_stmt    = IDENT "=" expr ";" ;
if_stmt        = "if" expr "{" { statement } "}"
                 "else" "{" { statement } "}" ;
for_range_stmt = "for" IDENT "in"
                 INT_LIT ".." INT_LIT
                 "{" { statement } "}" ;
swap_stmt      = "swap" "(" IDENT "," IDENT ")" ";" ;
incr_stmt      = "incr" "(" IDENT "," expr ")" ";" ;
decr_stmt      = "decr" "(" IDENT "," expr ")" ";" ;
xor_assign_stmt = IDENT "^=" expr ";" ;
trace_stmt     = "trace" "(" STRING_LIT
                 { "," expr } ")" ";" ;
checkpoint_stmt = "checkpoint" "(" STRING_LIT ")" ";" ;
assert_stmt    = "assert_invariant" "("
                 expr "," STRING_LIT ")" ";" ;
return_stmt    = "return" [ expr ] ";" ;
expr_stmt      = expr ";" ;
\end{lstlisting}

Note the absence of \texttt{while}, \texttt{loop}, and any production that
would permit a function body to reference its own name.  The \texttt{for\_range\_stmt}
production requires \texttt{INT\_LIT} (integer literal) for both bounds---not
\texttt{expr}---enforcing static bounds at the grammar level.


\section{Reversibility Proof Details}
\label{app:reversibility}

We provide the full reversibility proof for bounded loops, which is the most
involved case from \cref{thm:reversibility}.

\begin{lemma}[Bounded Loop Reversibility]
Let $L = \texttt{for}\;i\;\texttt{in}\;0\texttt{..}N\;\{s\}$ where $s$ is a
reversible statement.  Then $L$ is reversible, with inverse
$\inv{L} = \texttt{for}\;i\;\texttt{in}\;(N{-}1)\texttt{..}(-1)\;\{\inv{s}[i]\}$.
\end{lemma}

\begin{proof}
Let $\sigma_0$ be the initial state.  Forward execution produces the state
sequence:
\[
  \sigma_0 \xrightarrow{s[0]} \sigma_1 \xrightarrow{s[1]} \sigma_2
  \xrightarrow{s[2]} \cdots \xrightarrow{s[N{-}1]} \sigma_N
\]

The inverse loop executes $\inv{s}$ with indices $N{-}1, N{-}2, \ldots, 0$:
\[
  \sigma_N \xrightarrow{\inv{s}[N{-}1]} \sigma_{N-1}
  \xrightarrow{\inv{s}[N{-}2]} \sigma_{N-2}
  \xrightarrow{\inv{s}[N{-}3]} \cdots
  \xrightarrow{\inv{s}[0]} \sigma_0
\]

Each step $\inv{s}[k]$ restores $\sigma_k$ from $\sigma_{k+1}$ by the
assumption that $s$ is reversible.  The composition restores $\sigma_0$.

The trace produced by the forward loop has exactly $N$ entries; the trace
produced by the inverse loop appends $N$ additional entries (recording the
reversal).  Both traces are finite because $N$ is a compile-time constant.
\end{proof}


\section{Supported Cryptographic Algorithms}
\label{app:crypto}

\begin{table}[htbp]
\centering
\caption{Post-quantum and classical algorithms available via Zig FFI}
\begin{tabular}{@{}llll@{}}
\toprule
\textbf{Algorithm} & \textbf{Type} & \textbf{Standard} & \textbf{Library} \\
\midrule
ML-KEM-1024 (Kyber) & KEM & FIPS 203 & liboqs \\
ML-DSA-87 (Dilithium) & Signature & FIPS 204 & liboqs \\
SLH-DSA-256f (SPHINCS+) & Signature & FIPS 205 & liboqs \\
Ed25519 & Signature & RFC 8032 & libsodium \\
X25519 & Key exchange & RFC 7748 & libsodium \\
XChaCha20-Poly1305 & AEAD & RFC 8439 & libsodium \\
Argon2id & KDF & RFC 9106 & libsodium \\
\bottomrule
\end{tabular}
\end{table}


\end{document}
